% ============================================================
% §5.1 Q1结果分析 — 马尔可夫演化预测 (Stage 05 产物)
% 竞赛: 电工杯 B题 (2026)
% ============================================================

\section{Q1求解结果与分析}

\subsection{老年人口演化的马尔可夫递推结果}

基于§3.2构造的增生型转移矩阵$\mathbf{A}$，对10个小区独立运行5年递推预测。全区域三类老人总量的演化轨迹如图\ref{fig:q1_total}所示，各小区基年(t=0)与第5年末(t=5)的人口结构对比见图\ref{fig:q1_stacked}。

\begin{figure}[H]
\centering
\includegraphics[width=0.85\textwidth]{figures/figure_q1_total_trend.png}
\caption{全区域三类老人总量5年演化轨迹。自理老人(+5.7\%)与失能老人(+80.3\%)同步增长，半失能老人(-3.7\%)因净流出效应微降。}
\label{fig:q1_total}
\end{figure}

\begin{figure}[H]
\centering
\includegraphics[width=0.92\textwidth]{figures/figure_q1_stacked_bar.png}
\caption{基年与第5年末各小区老年人口结构对比。失能老人(橙色)占比从9.1\%跃升至14.8\%，构成"银发海啸"的底层推力。}
\label{fig:q1_stacked}
\end{figure}

表\ref{tab:q1_summary}汇总了全区域关键指标变化。

\begin{table}[H]
\centering
\caption{Q1人口预测关键指标汇总}
\label{tab:q1_summary}
\begin{tabular}{lrrr}
\toprule
\textbf{指标} & \textbf{基年(t=0)} & \textbf{第5年末(t=5)} & \textbf{变化(\%)} \\
\midrule
60+总人口 & 6~864 & 7~577 & +10.4\% \\
自理老人 & 4~712 & 4~981 & +5.7\% \\
半失能老人 & 1~528 & 1~471 & $-$3.7\% \\
失能老人 & 624 & 1~125 & \textbf{+80.3\%} \\
失能占总比 & 9.1\% & 14.8\% & +5.7pp \\
理论月需求总量 & -- & 262~553次/月 & -- \\
实际月需求总量 & -- & 247~060次/月 & -- \\
消费约束削减 & -- & 15~493次/月 & 5.9\% \\
\bottomrule
\end{tabular}
\end{table}

\subsubsection{核心发现与机制分析}

表\ref{tab:q1_summary}的核心发现是失能老人的\textbf{"失能棘轮效应"}：失能状态在转移矩阵$\mathbf{A}$中处于半吸收地位，对角元$a_{33}=0.95$，持续接收半失能群体10\%的年化流入（$a_{32}=0.095$），5年增幅达+80.3\%。半失能群体充当从自理到失能的加速通道，自身净减少3.7\%。每年7\%的新增自理注入产生总量稀释与结构滞后两重效应，使当前失能激增实际源自5年前已处于半失能状态的人群。

消费约束引入后，理论月需求削减5.9\%至247~060次/月，失能老人全10小区触发等比例削减（缺口23.9\%--49.1\%），构成Q3定价与补贴优化的底层经济学逻辑。Q1输出已结构化为两个CSV接口文件（\texttt{q1\_final\_population.csv}与\texttt{q1\_service\_demand.csv}）供Q2/Q3无缝读取，实现全流程数值闭环。
